\documentclass[10pt,a4paper]{article}

% ============================================================================
%  Sección Cocinas de TODO DECO — Optimización 2026 2C
%  Martínez Folica, Nazareno · Naón, Pedro
%
%  Compilar con dos pasadas de pdflatex (índice y referencias cruzadas).
%
%  Nota de maquetación: todas las tablas y figuras usan [H] (paquete float)
%  para que aparezcan exactamente donde están declaradas. El documento tiene
%  muchas tablas y con flotantes libres el orden de lectura se rompe: una
%  tabla puede subir al tope de la página y partir al medio una enumeración.
% ============================================================================

\usepackage[utf8]{inputenc}
\usepackage[T1]{fontenc}
\usepackage{lmodern}
\usepackage[a4paper,margin=2.2cm,headheight=14pt]{geometry}
\usepackage{graphicx}
\usepackage{booktabs}
\usepackage{array}
\usepackage{tabularx}
\usepackage{float}
\usepackage[table]{xcolor}
\usepackage{amsmath}
\usepackage{amssymb}
\usepackage{enumitem}
\usepackage{caption}
\captionsetup{font=small,skip=4pt}
\usepackage{fancyhdr}
\usepackage{titlesec}
\usepackage{tikz}
\usetikzlibrary{arrows.meta,calc,positioning,decorations.pathreplacing}
\usepackage[colorlinks=true,linkcolor=accentlink,citecolor=accentlink,urlcolor=accentlink,linktoc=all]{hyperref}

% ---------- paleta ----------------------------------------------------------
\definecolor{limita}{HTML}{2A78D6}
\definecolor{accent}{HTML}{1F4E79}
\definecolor{accentlink}{HTML}{1F5C99}

% ---------- tipografía de títulos -------------------------------------------
\titleformat{\section}{\normalfont\Large\bfseries\color{accent}}{\thesection}{0.6em}{}
\titleformat{\subsection}{\normalfont\large\bfseries\color{accent}}{\thesubsection}{0.6em}{}
\titleformat{\subsubsection}{\normalfont\bfseries\color{accent}}{\thesubsubsection}{0.6em}{}
\titlespacing*{\section}{0pt}{1.1em}{0.55em}
\titlespacing*{\subsection}{0pt}{0.9em}{0.4em}

\pagestyle{fancy}
\fancyhf{}
\renewcommand{\headrulewidth}{0.4pt}
\fancyhead[L]{\small Optimización 2026 2C — Sección Cocinas de TODO DECO}
\fancyhead[R]{\small Martínez Folica, N. · Naón, P.}
\fancyfoot[C]{\small\thepage}

\setlength{\parskip}{0.25em}
\setlist[itemize]{topsep=0.2em, itemsep=0.05em}
\setlist[enumerate]{topsep=0.2em, itemsep=0.05em}

\newcommand{\cuello}[1]{\textbf{\textcolor{limita}{#1}}}
\newcommand{\art}[1]{\texttt{#1}}
\newcommand{\cita}[1]{\emph{``#1''}}

\renewcommand{\contentsname}{Índice}
\renewcommand{\figurename}{Figura}
\renewcommand{\tablename}{Tabla}

\begin{document}

% ============================================================================
\begin{titlepage}
\begin{center}
\vspace*{1.2cm}
{\large UCA}\\[2pt]
{\large Facultad de Ingeniería y Ciencias Agrarias}\\[1.4cm]

{\large Optimización — 2026, 2\textdegree{} Cuatrimestre}\\[0.3cm]
{\large Trabajo Práctico N.\textdegree{} 2}\\[1.8cm]

\rule{0.85\linewidth}{0.6pt}\\[0.5cm]
{\Huge\bfseries\color{accent} Sección Cocinas de TODO DECO}\\[0.5cm]
{\Large Modelo de programación entera y binaria para la\\ definición del plan de stock del depósito}\\[0.5cm]
\rule{0.85\linewidth}{0.6pt}\\[2.2cm]

\begin{tabular}{c}
{\large\bfseries Integrantes} \\[0.4cm]
{\large Martínez Folica, Nazareno} \\[0.15cm]
{\large Naón, Pedro}
\end{tabular}

\vfill

{\large 17 de septiembre de 2026}
\end{center}
\end{titlepage}

\pagenumbering{roman}
\tableofcontents
\newpage
\pagenumbering{arabic}

% ============================================================================
\section{Situación y alcance}
\label{sec:situacion}

La cadena de artículos del hogar y construcción TODO DECO incorporó una Sección
Cocinas en una de sus sucursales del interior del país. La sección ofrece
\textbf{20 combos} de cocina, presentados comercialmente como combinaciones de
\cita{Cocina, lavavajillas, alacenas, mesadas, bacha, grifería, baldosas y papel
vinílico decorativo}. El catálogo que los abastece está organizado en
\textbf{ocho categorías} —baldosas, empapelado vinílico, apliques de luz,
alacenas, mesadas, bacha y grifería, lavavajillas y cocinas— que, entre todas
sus variantes, suman \textbf{30 artículos} distintos. Cada combo es una lista
cerrada de siete u ocho de esos artículos.

Conviene advertir que la enumeración comercial y el catálogo no coinciden: la
primera menciona por separado la bacha y la grifería, que se almacenan como una
sola categoría, y omite los apliques de luz, que sin embargo integran los veinte
combos. Este informe trabaja sobre las ocho categorías del catálogo.

La condición comercial que gobierna el problema está definida por la empresa en
una sola frase:

\begin{quote}
\cita{Los clientes pueden adquirir un combo completo a precio promocional si
todos los artículos se encuentran disponibles en stock en el momento del
pedido.}
\end{quote}

El espacio de almacenamiento, en cambio, es escaso. La empresa establece que
\cita{como el espacio es compartido con otras secciones y el almacén es pequeño,
la capacidad es muy limitada}, y fija un tope por tipo de artículo que totaliza
\textbf{33 lugares}. De esas dos definiciones surge la tensión que el modelo
tiene que resolver: \textbf{los combos comparten artículos y compiten por los
mismos lugares del depósito}, y basta con que falte uno solo de los artículos de
un combo para que ese combo deje de estar disponible para la venta.

El objetivo fijado para la sección es \textbf{maximizar la variedad de combos
disponibles para la venta}, y la decisión que hay que tomar para lograrlo es
cuántas unidades de cada uno de los 30 artículos mantener en stock.

\subsection{Alcance del análisis}

El informe formula un único modelo y lo resuelve bajo tres escenarios
sucesivos, que se corresponden con las tres situaciones que la empresa está
evaluando:

\begin{enumerate}
  \item \textbf{Reposición mensual, depósito actual} (Sección~\ref{sec:esc1}).
  Es la situación vigente: la reposición de los productos se realiza una vez al
  mes.
  \item \textbf{Reposición automática} (Sección~\ref{sec:esc2}). La empresa
  cerró \cita{un nuevo contrato con proveedores de la zona para que puedan
  disminuir considerablemente los tiempos de los envíos}, con lo cual la
  reposición pasa a considerarse instantánea.
  \item \textbf{Espacio compartido para mesadas y alacenas}
  (Sección~\ref{sec:esc3}). Una nueva disposición física del almacén permitiría
  reubicar ambas categorías en un espacio común de 12 unidades, sobre el
  escenario anterior.
\end{enumerate}

Los tres comparten formulación: cambian los parámetros, no las ecuaciones.

\subsection{Naturaleza del problema}

Las relaciones del problema son lineales, pero las cantidades involucradas no
son divisibles: no existe media cocina en stock ni medio combo disponible para
la venta. Las variables son, por lo tanto, \textbf{enteras y binarias}, y el
modelo pertenece a la programación lineal entera. Se resuelve mediante el
algoritmo de \emph{Branch \& Bound} \cite{catedra-mip, hillier}.

Esto tiene una consecuencia metodológica que conviene anticipar: en programación
entera \textbf{no se dispone del análisis post-óptimo} habitual de la
programación lineal continua \cite{catedra-mip}. El óptimo entero no varía de
forma continua con el término independiente de una restricción sino a saltos, de
modo que los precios sombra y los rangos permisibles pierden validez. El
análisis de sensibilidad se realiza, en su lugar, \textbf{re-resolviendo el
modelo}: se amplía la capacidad de un espacio en una unidad y se mide cuánto
varía el óptimo. Ese procedimiento se aplica en cada uno de los tres escenarios.

% ============================================================================
\section{Formulación del modelo}
\label{sec:modelo}

\subsection{Supuestos adoptados}
\label{sec:supuestos}

La información disponible es precisa en los conteos —la composición de cada
combo y la capacidad de cada espacio son datos exactos— pero deja abiertos
algunos puntos de interpretación que resultan indispensables para escribir el
modelo. Cada supuesto adoptado responde a tres preguntas: qué no puede
escribirse sin él, por qué se eligió esa lectura y no otra, y cómo se verificó
su incidencia sobre el resultado. Todos ellos son parámetros del modelo y no
constantes del código, de modo que evaluarlos consiste en resolver nuevamente
con otro valor.

\paragraph{S1 — La disponibilidad bajo reposición mensual se interpreta como
reserva exclusiva de unidades.} Es el supuesto que vincula el stock con los
combos, y sin él el modelo no puede formularse. Caben dos lecturas:

\begin{itemize}
  \item[(i)] \textbf{Exclusiva.} Cada combo ofrecido dispone de su propio juego
  de unidades reservado para el mes. Si dos combos incluyen \art{B2}, se
  requieren dos juegos de \art{B2}. Garantiza que cualquier combo ofrecido pueda
  venderse aun cuando ya se hayan vendido los demás.
  \item[(ii)] \textbf{Compartida.} Basta con una unidad de cada artículo: al
  inicio del mes todos los combos que la utilizan están disponibles, pero en
  cuanto se vende uno, los que comparten artículos con él dejan de estarlo hasta
  la reposición siguiente.
\end{itemize}

Se adopta la lectura \textbf{(i)}, por tres razones. La primera es que la
empresa precisa expresamente que \cita{la reposición de los productos se realiza
una vez al mes}, y bajo la lectura (ii) ese dato no tendría ninguna incidencia
sobre el modelo. La segunda es que la lectura (ii) describe exactamente la
situación de reposición automática analizada en la Sección~\ref{sec:esc2}:
adoptarla aquí volvería indistinguibles ambos escenarios. La tercera es
conceptual: la disponibilidad para la venta bajo reposición mensual debe
sostenerse \emph{durante} el mes, y no únicamente el primer día. La incidencia
de esta decisión se cuantifica en la propia Sección~\ref{sec:esc2}.

\paragraph{S2 — Cada combo incluye una unidad de cada artículo que lista.} La
composición de los combos indica \emph{qué} artículos los integran, no
\emph{cuántas} unidades de cada uno, mientras que las capacidades están
expresadas en juegos, rollos y unidades. Para baldosas, alacenas y bachas la
propia empresa habla de juegos, de modo que un juego por combo es la lectura
natural. El caso menos evidente es el empapelado, ya que una cocina real puede
requerir más de un rollo; por eso se evalúa un barrido de 1 a 9 rollos por combo
(Sección~\ref{sec:sens_s2}).

\paragraph{S3 — El desempate entre soluciones óptimas prioriza la cobertura del
catálogo y luego el ahorro de espacio.} El objetivo de variedad no incorpora el
stock, lo que deja indeterminados dos aspectos de la respuesta: qué conjunto de
combos se ofrece —hay cientos de conjuntos equivalentes— y cuánta mercadería se
guarda. Como la decisión buscada es precisamente el número de unidades por
artículo, se requiere un criterio que produzca una solución única. El adoptado
se apoya en dos definiciones de la propia empresa: que el espacio
\cita{es compartido con otras secciones}, de donde se sigue que a igual variedad
conviene el plan que menos lugar ocupa; y el catálogo de artículos que la
sección comercializa, cuyas ocho categorías no deberían quedar sin
representación en la oferta. El mecanismo se detalla en la
Sección~\ref{sec:objetivo} y se contrasta contra otros tres criterios en la
Sección~\ref{sec:sens_s3}.

\paragraph{S4 — Todas las variantes de una categoría ocupan el mismo lugar.} La
capacidad se mide en unidades sin distinguir variante: un juego de baldosas
\art{B1} ocupa lo mismo que uno \art{B4}. La empresa expresa las capacidades en
esos términos y lo confirma para el espacio compartido, al señalar que los
lavavajillas y las cocinas \cita{tienen el mismo tamaño}.

\paragraph{S5 — Todos los combos tienen igual valor y no se dispone de
pronósticos de demanda.} La variedad se define como la cantidad de combos
distintos disponibles, con peso unitario cada uno, ya que no se cuenta con
precios, márgenes ni estimaciones de venta. Ofrecer un combo significa poder
venderlo al menos una vez durante el mes, en coherencia con S1. De disponerse de
precios, el objetivo podría ponderarse sin modificar el resto del modelo.

\paragraph{S6 — Supuestos menores.} El horizonte de análisis es de un mes y el
stock inicial constituye el pico de ocupación, dado que luego solo desciende con
las ventas; la capacidad se controla, por lo tanto, contra ese stock inicial. El
depósito parte vacío, sin existencias previas ni pedidos pendientes. Los
artículos no se comercializan sueltos fuera de los combos, o bien esas ventas no
afectan la reserva. Rigen además los supuestos generales de todo modelo lineal
—proporcionalidad, aditividad, certidumbre y no negatividad—, con una salvedad
deliberada: la divisibilidad se abandona, y es precisamente lo que determina el
carácter entero del modelo.

\begin{table}[H]
\centering\small
\caption{Supuestos del modelo y alternativas evaluadas.}
\label{tab:supuestos}
\begin{tabularx}{\linewidth}{@{}llXl@{}}
\toprule
\textbf{ID} & \textbf{Supuesto} & \textbf{Lectura adoptada} & \textbf{Alternativas evaluadas} \\
\midrule
S1 & Disponibilidad         & Reserva exclusiva      & Reserva compartida \\
S2 & Unidades por artículo  & 1 por combo            & 1 a 9 rollos de empapelado \\
S3 & Desempate entre óptimos & Categorías y unidades & Otros tres criterios \\
S4 & Lugar por variante     & Uniforme               & --- \\
S5 & Valor de cada combo    & Unitario, sin demanda  & --- \\
S6 & Horizonte y stock inicial & Un mes, depósito vacío & --- \\
\bottomrule
\end{tabularx}
\end{table}

\subsection{Variables de decisión}
\label{sec:variables}

El modelo decide simultáneamente el stock a mantener y los combos que quedan
habilitados:
\[
x_p \in \mathbb{Z}_{\geq 0},\ p \in \{\art{B1}\dots\art{C4}\}
\qquad\qquad
y_c \in \{0,1\},\ c = 1\dots 20
\]
donde $x_p$ son las unidades del artículo $p$ en el depósito e $y_c$ vale 1 si
el combo $c$ queda disponible para la venta, es decir, si todos sus artículos
tienen unidad reservada. Son \textbf{50 variables}: 30 enteras y 20 binarias.

Mantener $x_p$ como variable explícita, en lugar de deducirla de los $y_c$,
responde a tres motivos: es la magnitud que la decisión requiere expresar; la
capacidad del depósito se formula naturalmente sobre el stock y no sobre los
combos; y en los escenarios siguientes lo que se modifica es la relación entre
stock y combos, de modo que basta reescribir la restricción de cobertura sin
alterar el resto de la formulación.

A estas dos familias se agrega una tercera, que interviene únicamente en las
etapas de desempate descriptas en la sección siguiente:
\[
z_k \in \{0,1\},\ k \in \{\text{baldosas}\dots\text{cocinas}\}
\]
con $z_k = 1$ si la categoría $k$ del catálogo queda representada en alguno de
los combos ofrecidos. Son 8 variables binarias adicionales.

Cabe una precisión técnica: bajo S1, declarar $x_p$ como entera es redundante,
puesto que si los $y_c$ son binarios el stock mínimo necesario resulta de una
suma de enteros. Se la declara entera de todos modos para que la formulación
refleje la naturaleza del problema \cite{catedra-mip} y no dependa de una
propiedad que dejaría de valer, por ejemplo, ante pedidos parciales.

\subsection{Función objetivo}
\label{sec:objetivo}

El objetivo de la sección se traduce de forma directa:
\[
\text{Max } V = \sum_{c=1}^{20} y_c
\]
con peso unitario para cada combo, según el supuesto S5.

Este funcional, sin embargo, resulta insuficiente para producir la decisión
buscada, por dos propiedades del problema que conviene explicitar:

\begin{enumerate}
  \item \textbf{El objetivo no incorpora el stock.} Guardar unidades que ningún
  combo utiliza no empeora el valor del funcional. Resolviendo únicamente esta
  etapa, la solución obtenida asigna \textbf{33 unidades} —la capacidad completa
  del depósito— para sostener tres combos que requieren 21: los 12 lugares
  restantes quedan ocupados con mercadería que no habilita ningún combo. La
  variedad sería correcta y el plan de stock, no.
  \item \textbf{La solución óptima no es única.} En el escenario base existen 854
  conjuntos de combos que alcanzan la variedad máxima
  (Sección~\ref{sec:esc1_alt}), y el algoritmo devuelve uno cualquiera de ellos.
\end{enumerate}

Ambas se resuelven incorporando dos objetivos subordinados, bajo el esquema de
\textbf{metas lexicográficas} \cite{catedra-metas}: cada etapa se optimiza sin
deteriorar el resultado de las anteriores.

\[
\begin{aligned}
\textbf{Etapa 1:}\quad & V^{*} = \text{Max} \sum_c y_c
  && \text{s.a. R1--R4} \\[2pt]
\textbf{Etapa 2:}\quad & K^{*} = \text{Max} \sum_k z_k
  && \text{s.a. R1--R4},\ \ \textstyle\sum_c y_c \geq V^{*} \\[2pt]
\textbf{Etapa 3:}\quad & U^{*} = \text{Min} \sum_p x_p
  && \text{s.a. R1--R4},\ \ \textstyle\sum_c y_c \geq V^{*},\ \sum_k z_k \geq K^{*}
\end{aligned}
\]

La etapa 1 determina cuántos combos pueden ofrecerse. La etapa 2 selecciona,
entre los planes que alcanzan esa variedad, aquellos que dejan representadas más
categorías del catálogo. La etapa 3 elige, entre estos últimos, el que menos
lugar ocupa del almacén compartido.

La etapa 2 merece justificación, porque atiende a una dificultad concreta del
problema. Aplicado en forma aislada, el criterio de mínimo stock favorece los
combos de siete artículos y admite soluciones que dejan una categoría entera del
catálogo sin representación: la terna 14--18--20, por ejemplo, alcanza la
variedad máxima con 21 unidades pero \textbf{no incluye ningún lavavajillas},
en contradicción con la descripción comercial de la sección. Sobre el diseño de
esta etapa se tomaron tres decisiones:

\begin{itemize}
  \item Se maximizan \textbf{las categorías del catálogo} en general, y no la
  presencia de un artículo puntual, de modo que el criterio conserve validez en
  los escenarios siguientes, donde la categoría ausente puede ser otra.
  \item Se incorpora como \textbf{objetivo y no como restricción}. Una
  restricción del tipo \emph{``al menos un combo con lavavajillas''} puede volver
  el modelo infactible, situación que efectivamente se verifica en la
  Sección~\ref{sec:sens_s2}. Como etapa de desempate, el modelo informa que
  cubre siete categorías en lugar de ocho y continúa resolviendo.
  \item La cobertura antecede al ahorro de espacio, porque la capacidad de
  ofrecer las categorías que la sección comercializa es condición de su
  propuesta comercial, mientras que liberar lugar beneficia a las demás
  secciones de la sucursal.
\end{itemize}

Una alternativa considerada fue un funcional único ponderado
$\text{Max} \sum_c y_c - \varepsilon \sum_p x_p$ con $\varepsilon < 1/33$, que
produce el mismo resultado porque ninguna reducción de stock compensa la pérdida
de un combo. Se optó por el esquema lexicográfico porque deja la jerarquía de
prioridades expresada en la formulación.

Por último, resolver el mismo modelo con variables continuas —$y_c \in [0,1]$ y
$x_p$ real— produce la \textbf{relajación lineal}, que aporta una cota superior
de la variedad y los precios sombra de cada espacio. Su alcance se discute en la
Sección~\ref{sec:esc1_relaj}.

\subsection{Parámetros}
\label{sec:parametros}

Los parámetros de entrada son tres: la composición de cada combo, las unidades
de cada artículo que un combo requiere y la capacidad de cada espacio del
depósito.
\[
a_{pc} \in \{0,1\}:\ \text{el combo } c \text{ incluye el artículo } p
\qquad
u_p:\ \text{unidades por combo}
\qquad
K_e:\ \text{capacidad del espacio } e
\]

\begin{table}[H]
\centering\small
\caption{Composición de los 20 combos.}
\label{tab:combos}
\begin{tabular}{rcccccccc r}
\toprule
\textbf{Combo} & Baldosa & Empap. & Aplique & Alacena & Mesada & Bacha & Lavav. & Cocina & \textbf{Art.} \\
\midrule
1  & B2 & E2 & L4 & A2 & M4 & G2 & W2 & C2 & 8 \\
2  & B1 & E1 & L1 & A4 & M4 & G4 & W1 & C2 & 8 \\
3  & B1 & E2 & L2 & A1 & M1 & G4 & W1 & C3 & 8 \\
4  & B3 & E3 & L3 & A3 & M3 & G1 & W1 & C1 & 8 \\
5  & B4 & E4 & L1 & A2 & M2 & G2 & W1 & C1 & 8 \\
6  & B2 & E2 & L2 & A4 & M4 & G3 & W2 & C4 & 8 \\
7  & B1 & E3 & L4 & A3 & M2 & G1 & W1 & C1 & 8 \\
8  & B2 & E1 & L3 & A1 & M1 & G3 & W2 & C4 & 8 \\
9  & B4 & E1 & L2 & A3 & M2 & G2 & W2 & C2 & 8 \\
10 & B1 & E4 & L1 & A1 & M3 & G4 & W1 & C3 & 8 \\
11 & B3 & E4 & L3 & A3 & M1 & G1 & W1 & C3 & 8 \\
12 & B2 & --- & L1 & A2 & M2 & G4 & W2 & C2 & 7 \\
13 & B4 & --- & L3 & A3 & M1 & G2 & W1 & C3 & 7 \\
14 & B4 & E1 & L4 & A1 & M3 & G1 & --- & C1 & 7 \\
15 & B3 & E2 & L1 & A1 & M1 & G3 & --- & C3 & 7 \\
16 & B3 & E3 & L4 & A1 & M3 & G2 & --- & C1 & 7 \\
17 & B1 & E3 & L2 & A3 & M3 & G4 & --- & C3 & 7 \\
18 & B2 & E3 & L3 & A2 & M4 & G1 & --- & C2 & 7 \\
19 & B2 & E2 & L4 & A4 & M4 & G2 & --- & C4 & 7 \\
20 & B2 & E2 & L1 & A1 & M2 & G3 & --- & C4 & 7 \\
\bottomrule
\end{tabular}
\vspace{1mm}\\
\footnotesize Once combos incluyen ocho artículos y nueve incluyen siete: los
combos 12 y 13 no llevan empapelado y los combos 14 a 20 no llevan lavavajillas.
\end{table}

Todos los combos, en cambio, incluyen exactamente una baldosa, un aplique, una
alacena, \textbf{una mesada}, una bacha y una cocina. Esa regularidad —en
particular la de las mesadas— determina el resultado del primer escenario.

\begin{table}[H]
\centering\small
\caption{Capacidad de almacenamiento por espacio.}
\label{tab:capacidad}
\begin{tabular}{llr}
\toprule
\textbf{Espacio} & \textbf{Artículos} & \textbf{Capacidad} \\
\midrule
Baldosas                & \art{B1}--\art{B4}                        & 5 juegos \\
Empapelado vinílico     & \art{E1}--\art{E4}                        & 8 rollos \\
Apliques de luz         & \art{L1}--\art{L4}                        & 4 unidades \\
Alacenas                & \art{A1}--\art{A4}                        & 4 juegos \\
Mesadas                 & \art{M1}--\art{M4}                        & 3 unidades \\
Bacha y grifería        & \art{G1}--\art{G4}                        & 4 juegos \\
Lavavajillas + cocinas  & \art{W1}--\art{W2}, \art{C1}--\art{C4}    & 5 unidades \\
\midrule
\textbf{Total} & & \textbf{33 lugares} \\
\bottomrule
\end{tabular}
\end{table}

La última fila de la Tabla~\ref{tab:capacidad} corresponde a una definición
explícita de la empresa: los lavavajillas y las cocinas \cita{tienen el mismo
tamaño así que pueden almacenarse juntos, con un máximo de 5 unidades en
total}. Es el único espacio que aloja dos categorías, y adelanta la forma que
tomará la reubicación evaluada en la Sección~\ref{sec:esc3}.

\begin{table}[H]
\centering\small
\caption{Cantidad de combos en los que interviene cada artículo.}
\label{tab:apariciones}
\begin{tabular}{lcccc r}
\toprule
\textbf{Categoría} & \textbf{Var. 1} & \textbf{Var. 2} & \textbf{Var. 3} & \textbf{Var. 4} & \textbf{Total} \\
\midrule
Baldosas            & B1: 5 & B2: 7 & B3: 4 & B4: 4 & 20 \\
Empapelado vinílico & E1: 4 & E2: 6 & E3: 5 & E4: 3 & 18 \\
Apliques de luz     & L1: 6 & L2: 4 & L3: 5 & L4: 5 & 20 \\
Alacenas            & A1: 7 & A2: 4 & A3: 6 & A4: 3 & 20 \\
Mesadas             & M1: 5 & M2: 5 & M3: 5 & M4: 5 & 20 \\
Bacha y grifería    & G1: 5 & G2: 6 & G3: 4 & G4: 5 & 20 \\
Lavavajillas        & W1: 8 & W2: 5 & ---   & ---   & 13 \\
Cocinas             & C1: 5 & C2: 5 & C3: 6 & C4: 4 & 20 \\
\midrule
\textbf{Total} & & & & & \textbf{151} \\
\bottomrule
\end{tabular}
\vspace{1mm}\\
\footnotesize Control: $11 \times 8 + 9 \times 7 = 151$ artículos requeridos en
total.
\end{table}

Los valores de la Tabla~\ref{tab:apariciones} son los coeficientes de la
restricción de cobertura bajo reserva exclusiva, y permiten dimensionar el
problema antes de resolverlo: reservar \art{B2} para los siete combos que lo
utilizan exigiría siete juegos de baldosas, frente a los cinco que admite el
depósito.

\begin{table}[H]
\centering\small
\caption{Cota máxima de variedad que impone cada espacio, bajo reserva
exclusiva. Si todos los combos ocupan al menos $u$ lugares de un espacio de
capacidad $K$, no pueden ofrecerse más de $\lfloor K/u \rfloor$ combos.}
\label{tab:cota_trivial}
\begin{tabular}{lrrr}
\toprule
\textbf{Espacio} & \textbf{Capacidad} & \textbf{Lugares por combo} & \textbf{Cota} \\
\midrule
Baldosas               & 5 & 1 & 5 \\
Empapelado vinílico    & 8 & 0 & sin cota \\
Apliques de luz        & 4 & 1 & 4 \\
Alacenas               & 4 & 1 & 4 \\
\cuello{Mesadas}       & \cuello{3} & \cuello{1} & \cuello{3} \\
Bacha y grifería       & 4 & 1 & 4 \\
Lavavajillas + cocinas & 5 & 1 & 5 \\
\bottomrule
\end{tabular}
\end{table}

De la Tabla~\ref{tab:cota_trivial} se desprenden tres conclusiones anticipadas,
todas bajo la lectura exclusiva de S1:

\begin{enumerate}
  \item \textbf{Las mesadas fijan el techo del problema.} Todos los combos
  llevan exactamente una mesada y el depósito admite tres, de modo que no pueden
  ofrecerse más de tres combos. La optimización se reduce a determinar cuáles.
  \item \textbf{El espacio de lavavajillas y cocinas actúa como segundo
  filtro.} Tres combos tomados entre el 1 y el 13 ocupan $2+2+2 = 6$ lugares y
  el espacio admite 5; por lo tanto, de cada tres combos ofrecidos a lo sumo dos
  pueden pertenecer a ese grupo.
  \item \textbf{El empapelado no restringe en ningún caso.} Con tres combos se
  utilizan como máximo tres rollos de los ocho disponibles. Lo mismo ocurre con
  baldosas, apliques, alacenas y bachas, a los que les resta al menos un lugar.
\end{enumerate}

\subsection{Restricciones}
\label{sec:restricciones}

Las restricciones traducen las dos definiciones operativas de la empresa: la
condición de venta del combo completo y la capacidad del almacén.

La \textbf{restricción de cobertura (R1)} surge de la condición comercial citada
en la Sección~\ref{sec:situacion} combinada con la periodicidad mensual de la
reposición: si cada combo ofrecido requiere su propio juego de unidades, el
stock de un artículo debe alcanzar para todos los combos que lo incluyen.

La \textbf{restricción de capacidad (R2)} recoge la definición del depósito:
\cita{se pueden almacenar hasta 5 juegos de baldosas, 8 rollos de empapelado
vinílico, 4 apliques de luz, 4 juegos de alacenas, 3 mesadas y 4 juegos de bacha
y grifería}, con el agregado del espacio compartido de lavavajillas y cocinas,
que se formula como R2$'$.

\begin{table}[H]
\centering\small
\caption{Restricciones del modelo.}
\label{tab:restricciones}
\begin{tabular}{llrr}
\toprule
\textbf{\#} & \textbf{Concepto} & \textbf{Expresión} & \textbf{Cantidad} \\
\midrule
R1     & Cobertura de combos     & $x_p \geq \sum_c u_p\, a_{pc}\, y_c \quad \forall p$ & 30 \\
R2     & Capacidad por espacio   & $\sum_{p \in P_e} x_p \leq K_e$                      & 6 \\
R2$'$  & Espacio compartido      & $x_{W1}+x_{W2}+\sum_{j} x_{Cj} \leq 5$               & 1 \\
R3     & Integralidad            & $y_c \in \{0,1\}$, \ $x_p \in \mathbb{Z}$            & --- \\
R4     & No negatividad          & $x_p \geq 0 \quad \forall p$                         & --- \\
\bottomrule
\end{tabular}
\vspace{1mm}\\
\footnotesize Las etapas 2 y 3 del desempate agregan 8 variables binarias $z_k$
y sus 8 restricciones de atadura $z_k \leq \sum_{c \,:\, k \in c} y_c$.
\end{table}

El modelo queda conformado por \textbf{37 restricciones} en la etapa 1 y 45 en
su versión completa, sobre 58 variables. R2 y R2$'$ son, en rigor, la misma
restricción aplicada a distintos conjuntos de artículos $P_e$: el depósito se
representa como una lista de espacios, cada uno con las categorías que aloja y
su capacidad. Esa formulación es la que permite evaluar la reubicación de la
Sección~\ref{sec:esc3} modificando un dato y no el modelo.

Dos ejemplos escritos en forma completa:
\[
\text{R1 para \art{B2}:}\quad x_{B2} \geq y_{1}+y_{6}+y_{8}+y_{12}+y_{18}+y_{19}+y_{20}
\qquad\qquad
\text{R2 mesadas:}\quad \sum_{j=1}^{4} x_{Mj} \leq 3
\]

\paragraph{Cota analítica de la variedad.} Sumando R1 sobre las cuatro variantes
de mesada, y dado que cada combo incluye exactamente una —de modo que cada $y_c$
interviene una sola vez en esa suma—, se obtiene:
\[
\sum_{j=1}^{4} x_{Mj} \;\geq\; \sum_{c} y_c
\qquad\text{y por R2}\qquad
\sum_{j=1}^{4} x_{Mj} \;\leq\; 3
\qquad\Longrightarrow\qquad
\boxed{\ \sum_c y_c \leq 3\ }
\]

La cota es alcanzable: cualquier terna tomada entre los combos 14 y 20 ocupa
exactamente tres lugares en cada uno de los siete espacios, todos con capacidad
igual o superior a tres. La variedad máxima del escenario base queda así
determinada analíticamente en \textbf{tres combos}, con independencia del
solver, lo que constituye el primer control de validez del modelo.

\subsection{Esquema del proceso}
\label{sec:esquema}

La Figura~\ref{fig:proceso} representa el problema en su conjunto: la capacidad
de cada espacio del depósito limita el stock disponible, el stock determina qué
combos quedan completos y la cantidad de combos completos constituye la variedad
que se maximiza.

\begin{figure}[H]
\centering
\resizebox{\linewidth}{!}{% =====================================================================
% Esquema del proceso — TP2, Sección Cocinas de TODO DECO (punto a).
% Es la versión prolija del diagrama de docs/plan_de_trabajo.md §9.
%
% Uso en el informe:
%     \usepackage{tikz}
%     \usetikzlibrary{positioning, arrows.meta, calc, decorations.pathreplacing}
%     ...
%     \begin{figure}[htbp]
%       \centering
%       % =====================================================================
% Esquema del proceso — TP2, Sección Cocinas de TODO DECO (punto a).
% Es la versión prolija del diagrama de docs/plan_de_trabajo.md §9.
%
% Uso en el informe:
%     \usepackage{tikz}
%     \usetikzlibrary{positioning, arrows.meta, calc, decorations.pathreplacing}
%     ...
%     \begin{figure}[htbp]
%       \centering
%       % =====================================================================
% Esquema del proceso — TP2, Sección Cocinas de TODO DECO (punto a).
% Es la versión prolija del diagrama de docs/plan_de_trabajo.md §9.
%
% Uso en el informe:
%     \usepackage{tikz}
%     \usetikzlibrary{positioning, arrows.meta, calc, decorations.pathreplacing}
%     ...
%     \begin{figure}[htbp]
%       \centering
%       \input{../docs/esquema_proceso.tex}
%       \caption{Del stock del depósito a la variedad de combos ofrecidos.}
%     \end{figure}
%
% Los colores son los mismos de los gráficos (src/graficos.py).
%
% Nota: el esquema no usa los caracteres "<" ni ">" en ningún lado (las flechas se
% declaran como -{Stealth} y las desigualdades con \leq). Es a propósito: babel en
% español los toma como atajos de «comillas» y rompería la compilación.
% =====================================================================

\definecolor{tpAcento}{HTML}{2A78D6}
\definecolor{tpAcentoSuave}{HTML}{CDE2FB}
\definecolor{tpGris}{HTML}{898781}
\definecolor{tpGrisSuave}{HTML}{E1E0D9}
\definecolor{tpTinta}{HTML}{0B0B0B}
\definecolor{tpTintaDos}{HTML}{52514E}

\begin{tikzpicture}[
    font=\footnotesize,
    espacio/.style={
      draw=tpGrisSuave, fill=tpGrisSuave!35, rounded corners=2pt,
      text=tpTinta, minimum width=4.9cm, minimum height=0.48cm,
      inner sep=3pt, align=center,
    },
    limita/.style={espacio, draw=tpAcento, fill=tpAcentoSuave},
    caja/.style={
      draw=tpGris, rounded corners=3pt, text=tpTinta,
      minimum width=3.1cm, minimum height=1.5cm, inner sep=5pt, align=center,
    },
    resultado/.style={caja, draw=tpAcento, fill=tpAcentoSuave},
    flecha/.style={-{Stealth[length=4pt]}, draw=tpGris, line width=0.7pt},
    etiqueta/.style={text=tpTintaDos, font=\scriptsize, midway, above},
  ]

  % --- Columna 1: los espacios del depósito (recurso) ---------------
  \node[espacio] (e1) {Baldosas \quad $\leq 5$};
  \node[espacio, below=1.6mm of e1] (e2) {Empapelado \quad $\leq 8$};
  \node[espacio, below=1.6mm of e2] (e3) {Apliques de luz \quad $\leq 4$};
  \node[espacio, below=1.6mm of e3] (e4) {Alacenas \quad $\leq 4$};
  \node[limita,  below=1.6mm of e4] (e5) {\textbf{Mesadas} \quad $\boldsymbol{\leq 3}$};
  \node[espacio, below=1.6mm of e5] (e6) {Bacha y grifería \quad $\leq 4$};
  \node[espacio, below=1.6mm of e6] (e7) {Lavavajillas + cocinas \quad $\leq 5$};

  \node[text=tpTintaDos, font=\scriptsize\bfseries, above=2.5mm of e1] (tituloDep)
    {DEPÓSITO (capacidad por espacio)};

  % Llave que agrupa los siete espacios.
  \coordinate (depTop) at ($(e1.north east)+(2mm,0)$);
  \coordinate (depBot) at ($(e7.south east)+(2mm,0)$);
  \draw[draw=tpGris, decorate, decoration={brace, amplitude=4pt}]
    (depTop) -- (depBot);

  % --- Columna 2: el stock (variables de decisión) ------------------
  \node[caja, right=15mm of e4] (stock)
    {\textbf{Stock}\\[1mm] $x_p \in \mathbb{Z}_{\geq 0}$\\[0.5mm] unidades de cada\\ uno de los 30 artículos};

  % --- Columna 3: los combos ----------------------------------------
  \node[caja, right=15mm of stock] (combos)
    {\textbf{Combos completos}\\[1mm] $y_c \in \{0,1\}$\\[0.5mm] 20 combos\\ de 7 u 8 artículos};

  % --- Columna 4: el resultado --------------------------------------
  \node[resultado, right=15mm of combos] (variedad)
    {\textbf{Variedad}\\[1mm] $\max \sum_c y_c$\\[0.5mm] combos disponibles\\ para la venta};

  % --- Flechas -------------------------------------------------------
  \draw[flecha] ($(depTop)!0.5!(depBot)+(4mm,0)$) -- (stock.west)
    node[etiqueta, align=center] {R2\\ capacidad};
  \draw[flecha] (stock.east) -- (combos.west)
    node[etiqueta, align=center] {R1\\ cobertura};
  \draw[flecha] (combos.east) -- (variedad.west);

  % --- Nota al pie del esquema ---------------------------------------
  \node[text=tpTintaDos, font=\scriptsize, align=center, below=6mm of combos] (nota)
    {Un combo cuenta solo si \emph{todos} sus artículos tienen\\
     una unidad reservada para él (supuesto S1).};

\end{tikzpicture}

%       \caption{Del stock del depósito a la variedad de combos ofrecidos.}
%     \end{figure}
%
% Los colores son los mismos de los gráficos (src/graficos.py).
%
% Nota: el esquema no usa los caracteres "<" ni ">" en ningún lado (las flechas se
% declaran como -{Stealth} y las desigualdades con \leq). Es a propósito: babel en
% español los toma como atajos de «comillas» y rompería la compilación.
% =====================================================================

\definecolor{tpAcento}{HTML}{2A78D6}
\definecolor{tpAcentoSuave}{HTML}{CDE2FB}
\definecolor{tpGris}{HTML}{898781}
\definecolor{tpGrisSuave}{HTML}{E1E0D9}
\definecolor{tpTinta}{HTML}{0B0B0B}
\definecolor{tpTintaDos}{HTML}{52514E}

\begin{tikzpicture}[
    font=\footnotesize,
    espacio/.style={
      draw=tpGrisSuave, fill=tpGrisSuave!35, rounded corners=2pt,
      text=tpTinta, minimum width=4.9cm, minimum height=0.48cm,
      inner sep=3pt, align=center,
    },
    limita/.style={espacio, draw=tpAcento, fill=tpAcentoSuave},
    caja/.style={
      draw=tpGris, rounded corners=3pt, text=tpTinta,
      minimum width=3.1cm, minimum height=1.5cm, inner sep=5pt, align=center,
    },
    resultado/.style={caja, draw=tpAcento, fill=tpAcentoSuave},
    flecha/.style={-{Stealth[length=4pt]}, draw=tpGris, line width=0.7pt},
    etiqueta/.style={text=tpTintaDos, font=\scriptsize, midway, above},
  ]

  % --- Columna 1: los espacios del depósito (recurso) ---------------
  \node[espacio] (e1) {Baldosas \quad $\leq 5$};
  \node[espacio, below=1.6mm of e1] (e2) {Empapelado \quad $\leq 8$};
  \node[espacio, below=1.6mm of e2] (e3) {Apliques de luz \quad $\leq 4$};
  \node[espacio, below=1.6mm of e3] (e4) {Alacenas \quad $\leq 4$};
  \node[limita,  below=1.6mm of e4] (e5) {\textbf{Mesadas} \quad $\boldsymbol{\leq 3}$};
  \node[espacio, below=1.6mm of e5] (e6) {Bacha y grifería \quad $\leq 4$};
  \node[espacio, below=1.6mm of e6] (e7) {Lavavajillas + cocinas \quad $\leq 5$};

  \node[text=tpTintaDos, font=\scriptsize\bfseries, above=2.5mm of e1] (tituloDep)
    {DEPÓSITO (capacidad por espacio)};

  % Llave que agrupa los siete espacios.
  \coordinate (depTop) at ($(e1.north east)+(2mm,0)$);
  \coordinate (depBot) at ($(e7.south east)+(2mm,0)$);
  \draw[draw=tpGris, decorate, decoration={brace, amplitude=4pt}]
    (depTop) -- (depBot);

  % --- Columna 2: el stock (variables de decisión) ------------------
  \node[caja, right=15mm of e4] (stock)
    {\textbf{Stock}\\[1mm] $x_p \in \mathbb{Z}_{\geq 0}$\\[0.5mm] unidades de cada\\ uno de los 30 artículos};

  % --- Columna 3: los combos ----------------------------------------
  \node[caja, right=15mm of stock] (combos)
    {\textbf{Combos completos}\\[1mm] $y_c \in \{0,1\}$\\[0.5mm] 20 combos\\ de 7 u 8 artículos};

  % --- Columna 4: el resultado --------------------------------------
  \node[resultado, right=15mm of combos] (variedad)
    {\textbf{Variedad}\\[1mm] $\max \sum_c y_c$\\[0.5mm] combos disponibles\\ para la venta};

  % --- Flechas -------------------------------------------------------
  \draw[flecha] ($(depTop)!0.5!(depBot)+(4mm,0)$) -- (stock.west)
    node[etiqueta, align=center] {R2\\ capacidad};
  \draw[flecha] (stock.east) -- (combos.west)
    node[etiqueta, align=center] {R1\\ cobertura};
  \draw[flecha] (combos.east) -- (variedad.west);

  % --- Nota al pie del esquema ---------------------------------------
  \node[text=tpTintaDos, font=\scriptsize, align=center, below=6mm of combos] (nota)
    {Un combo cuenta solo si \emph{todos} sus artículos tienen\\
     una unidad reservada para él (supuesto S1).};

\end{tikzpicture}

%       \caption{Del stock del depósito a la variedad de combos ofrecidos.}
%     \end{figure}
%
% Los colores son los mismos de los gráficos (src/graficos.py).
%
% Nota: el esquema no usa los caracteres "<" ni ">" en ningún lado (las flechas se
% declaran como -{Stealth} y las desigualdades con \leq). Es a propósito: babel en
% español los toma como atajos de «comillas» y rompería la compilación.
% =====================================================================

\definecolor{tpAcento}{HTML}{2A78D6}
\definecolor{tpAcentoSuave}{HTML}{CDE2FB}
\definecolor{tpGris}{HTML}{898781}
\definecolor{tpGrisSuave}{HTML}{E1E0D9}
\definecolor{tpTinta}{HTML}{0B0B0B}
\definecolor{tpTintaDos}{HTML}{52514E}

\begin{tikzpicture}[
    font=\footnotesize,
    espacio/.style={
      draw=tpGrisSuave, fill=tpGrisSuave!35, rounded corners=2pt,
      text=tpTinta, minimum width=4.9cm, minimum height=0.48cm,
      inner sep=3pt, align=center,
    },
    limita/.style={espacio, draw=tpAcento, fill=tpAcentoSuave},
    caja/.style={
      draw=tpGris, rounded corners=3pt, text=tpTinta,
      minimum width=3.1cm, minimum height=1.5cm, inner sep=5pt, align=center,
    },
    resultado/.style={caja, draw=tpAcento, fill=tpAcentoSuave},
    flecha/.style={-{Stealth[length=4pt]}, draw=tpGris, line width=0.7pt},
    etiqueta/.style={text=tpTintaDos, font=\scriptsize, midway, above},
  ]

  % --- Columna 1: los espacios del depósito (recurso) ---------------
  \node[espacio] (e1) {Baldosas \quad $\leq 5$};
  \node[espacio, below=1.6mm of e1] (e2) {Empapelado \quad $\leq 8$};
  \node[espacio, below=1.6mm of e2] (e3) {Apliques de luz \quad $\leq 4$};
  \node[espacio, below=1.6mm of e3] (e4) {Alacenas \quad $\leq 4$};
  \node[limita,  below=1.6mm of e4] (e5) {\textbf{Mesadas} \quad $\boldsymbol{\leq 3}$};
  \node[espacio, below=1.6mm of e5] (e6) {Bacha y grifería \quad $\leq 4$};
  \node[espacio, below=1.6mm of e6] (e7) {Lavavajillas + cocinas \quad $\leq 5$};

  \node[text=tpTintaDos, font=\scriptsize\bfseries, above=2.5mm of e1] (tituloDep)
    {DEPÓSITO (capacidad por espacio)};

  % Llave que agrupa los siete espacios.
  \coordinate (depTop) at ($(e1.north east)+(2mm,0)$);
  \coordinate (depBot) at ($(e7.south east)+(2mm,0)$);
  \draw[draw=tpGris, decorate, decoration={brace, amplitude=4pt}]
    (depTop) -- (depBot);

  % --- Columna 2: el stock (variables de decisión) ------------------
  \node[caja, right=15mm of e4] (stock)
    {\textbf{Stock}\\[1mm] $x_p \in \mathbb{Z}_{\geq 0}$\\[0.5mm] unidades de cada\\ uno de los 30 artículos};

  % --- Columna 3: los combos ----------------------------------------
  \node[caja, right=15mm of stock] (combos)
    {\textbf{Combos completos}\\[1mm] $y_c \in \{0,1\}$\\[0.5mm] 20 combos\\ de 7 u 8 artículos};

  % --- Columna 4: el resultado --------------------------------------
  \node[resultado, right=15mm of combos] (variedad)
    {\textbf{Variedad}\\[1mm] $\max \sum_c y_c$\\[0.5mm] combos disponibles\\ para la venta};

  % --- Flechas -------------------------------------------------------
  \draw[flecha] ($(depTop)!0.5!(depBot)+(4mm,0)$) -- (stock.west)
    node[etiqueta, align=center] {R2\\ capacidad};
  \draw[flecha] (stock.east) -- (combos.west)
    node[etiqueta, align=center] {R1\\ cobertura};
  \draw[flecha] (combos.east) -- (variedad.west);

  % --- Nota al pie del esquema ---------------------------------------
  \node[text=tpTintaDos, font=\scriptsize, align=center, below=6mm of combos] (nota)
    {Un combo cuenta solo si \emph{todos} sus artículos tienen\\
     una unidad reservada para él (supuesto S1).};

\end{tikzpicture}
}
\caption{Del depósito a la oferta: la capacidad de cada espacio se transforma en
variedad de combos disponibles. En azul, el espacio que limita en el escenario
base. Las dos flechas corresponden a las dos familias de restricciones.}
\label{fig:proceso}
\end{figure}

El esquema corresponde al escenario base. Los dos escenarios siguientes no
modifican su estructura sino dos elementos dentro de ella: la reposición
automática altera la relación de cobertura, y la nueva disposición física
fusiona dos espacios de la columna izquierda en uno solo.

% ============================================================================
\section{Escenario 1 — Reposición mensual}
\label{sec:esc1}

Corresponde a la situación vigente, en la que la reposición de los productos se
realiza una vez al mes y el depósito conserva su disposición actual. El modelo
se resolvió con el solver CBC a través de PuLP, con estado \texttt{Optimal} en
las tres etapas (véase el Anexo).

\subsection{Variedad alcanzada y plan de stock}

La variedad máxima es de \textbf{3 combos sobre 20}, el 15\,\% del catálogo, en
coincidencia con la cota analítica deducida en la
Sección~\ref{sec:restricciones}. El plan de stock resultante requiere
\textbf{21 unidades} distribuidas en \textbf{17 de los 30 artículos}, y habilita
los combos \textbf{12, 15 y 18}.

\begin{table}[H]
\centering\small
\caption{Combos habilitados y su requerimiento de stock.}
\label{tab:combos_e1}
\begin{tabular}{rlrr}
\toprule
\textbf{Combo} & \textbf{Artículos} & \textbf{Unidades} & \textbf{Lugares lavav. + cocinas} \\
\midrule
12 & \art{B2 L1 A2 M2 G4 W2 C2} & 7 & 2 \\
15 & \art{B3 E2 L1 A1 M1 G3 C3} & 7 & 1 \\
18 & \art{B2 E3 L3 A2 M4 G1 C2} & 7 & 1 \\
\midrule
\textbf{Total} & & \textbf{21} & \textbf{4} \\
\bottomrule
\end{tabular}
\end{table}

La combinación cubre las ocho categorías del catálogo: el combo 12 aporta el
lavavajillas y los combos 15 y 18 aportan el empapelado, categorías que ninguno
de ellos incluye simultáneamente.

\begin{table}[H]
\centering\footnotesize
\renewcommand{\arraystretch}{0.95}
\caption{Plan de stock. Los trece artículos restantes quedan en cero: \art{B1},
\art{B4}, \art{E1}, \art{E4}, \art{L2}, \art{L4}, \art{A3}, \art{A4}, \art{M3},
\art{G2}, \art{W1}, \art{C1} y \art{C4}.}
\label{tab:stock_e1}
\begin{tabular}{llrl}
\toprule
\textbf{Categoría} & \textbf{Artículo} & \textbf{Unidades} & \textbf{Combos} \\
\midrule
Baldosas            & \art{B2} Marfil                        & 2 & 12, 18 \\
Baldosas            & \art{B3} Blanco y azul a cuadros       & 1 & 15 \\
Empapelado vinílico & \art{E2} Símil marfil a rayas celestes & 1 & 15 \\
Empapelado vinílico & \art{E3} Símil mármol azul             & 1 & 18 \\
Apliques de luz     & \art{L1} Plafón único rectangular      & 2 & 12, 15 \\
Apliques de luz     & \art{L3} Bombillas de filamentos       & 1 & 18 \\
Alacenas            & \art{A1} Madera clara                  & 1 & 15 \\
Alacenas            & \art{A2} Madera oscura                 & 2 & 12, 18 \\
Mesadas             & \art{M1} Madera laqueada               & 1 & 15 \\
Mesadas             & \art{M2} Cemento alisado               & 1 & 12 \\
Mesadas             & \art{M4} Granito                       & 1 & 18 \\
Bacha y grifería    & \art{G1} Bacha dividida, grifo mono-comando & 1 & 18 \\
Bacha y grifería    & \art{G3} Bacha única, grifo mono-comando    & 1 & 15 \\
Bacha y grifería    & \art{G4} Bacha única, grifos separados      & 1 & 12 \\
Lavavajillas        & \art{W2} Gris                          & 1 & 12 \\
Cocinas             & \art{C2} Cocina eléctrica negra        & 2 & 12, 18 \\
Cocinas             & \art{C3} Cocina a gas blanca           & 1 & 15 \\
\midrule
\textbf{Total} & \textbf{17 artículos} & \textbf{21} & \\
\bottomrule
\end{tabular}
\end{table}

Cuatro artículos requieren dos unidades —\art{B2}, \art{L1}, \art{A2} y
\art{C2}—: son precisamente los que comparten dos de los combos habilitados, y
constituyen la manifestación directa de la reserva exclusiva del supuesto S1.
Esa duplicación es la que desaparece al pasar a reposición automática.

\subsection{Ocupación del depósito}

\begin{table}[H]
\centering\small
\caption{Ocupación del depósito bajo el plan propuesto.}
\label{tab:ocupacion_e1}
\begin{tabular}{lrrr}
\toprule
\textbf{Espacio} & \textbf{Capacidad} & \textbf{Ocupado} & \textbf{Holgura} \\
\midrule
Baldosas               & 5 & 3 & 2 \\
Empapelado vinílico    & 8 & 2 & 6 \\
Apliques de luz        & 4 & 3 & 1 \\
Alacenas               & 4 & 3 & 1 \\
\cuello{Mesadas}       & \cuello{3} & \cuello{3} & \cuello{0} \\
Bacha y grifería       & 4 & 3 & 1 \\
Lavavajillas + cocinas & 5 & 4 & 1 \\
\midrule
\textbf{Total} & \textbf{33} & \textbf{21} & \textbf{12} \\
\bottomrule
\end{tabular}
\end{table}

Un solo espacio se agota, y lo hace en todas las soluciones óptimas: las mesadas
están completas en las 854 ternas que alcanzan la variedad máxima. Al mismo
tiempo, \textbf{12 de los 33 lugares del depósito —el 36\,\%— quedan sin
utilizar}. La limitación de la sección no proviene, por lo tanto, de la
dimensión general del almacén sino de un único espacio de tres lugares, lo que
modifica por completo el objeto de una eventual negociación de capacidad con las
demás secciones de la sucursal.

\subsection{Valor de una ampliación de capacidad}
\label{sec:esc1_cap}

Dado que el modelo entero no admite precios sombra, el valor de una unidad
adicional de almacenamiento se determina ampliando la capacidad de cada espacio
y resolviendo nuevamente.

\begin{table}[H]
\centering\small
\caption{Efecto de un lugar adicional, un espacio por vez.}
\label{tab:cap_e1}
\begin{tabular}{lrrr}
\toprule
\textbf{Espacio} & \textbf{Capacidad ampliada} & $\boldsymbol{V^{*}}$ & $\boldsymbol{\Delta V^{*}}$ \\
\midrule
Baldosas               & 6 & 3 & 0 \\
Empapelado vinílico    & 9 & 3 & 0 \\
Apliques de luz        & 5 & 3 & 0 \\
Alacenas               & 5 & 3 & 0 \\
\cuello{Mesadas}       & \cuello{4} & \cuello{4} & \cuello{$+1$} \\
Bacha y grifería       & 5 & 3 & 0 \\
Lavavajillas + cocinas & 6 & 3 & 0 \\
\bottomrule
\end{tabular}
\end{table}

\begin{table}[H]
\centering\small
\caption{Variedad máxima según la capacidad del espacio de mesadas.}
\label{tab:barrido_e1}
\begin{tabular}{rr}
\toprule
\textbf{Capacidad de mesadas} & $\boldsymbol{V^{*}}$ \\
\midrule
3 (actual) & 3 \\
4          & 4 \\
5          & 4 \\
6          & 4 \\
\bottomrule
\end{tabular}
\end{table}

Un lugar adicional de mesada habilita un combo; un lugar adicional en cualquier
otro espacio no produce ningún efecto. El barrido de la
Tabla~\ref{tab:barrido_e1} muestra además que el beneficio se agota de
inmediato: con cuatro mesadas la variedad queda limitada simultáneamente por
apliques, alacenas y bachas —los tres espacios tienen cuatro lugares y todos los
combos utilizan uno de cada uno—, de modo que a partir de allí ninguna
ampliación aislada incrementa la oferta. Para superar los cuatro combos sería
necesario ampliar varios espacios de forma conjunta.

\subsection{Relajación lineal}
\label{sec:esc1_relaj}

\begin{table}[H]
\centering\small
\caption{Relajación lineal del modelo. La cota obtenida es de 3,0000 combos y
la solución resulta entera.}
\label{tab:relajacion_e1}
\begin{tabular}{lrrr r}
\toprule
\textbf{Espacio} & \textbf{Capacidad} & \textbf{Ocupado} & \textbf{Holgura} & \textbf{Precio sombra} \\
\midrule
Baldosas               & 5 & 3 & 2 & 0 \\
Empapelado vinílico    & 8 & 3 & 5 & 0 \\
Apliques de luz        & 4 & 3 & 1 & 0 \\
Alacenas               & 4 & 3 & 1 & 0 \\
\cuello{Mesadas}       & \cuello{3} & \cuello{3} & \cuello{0} & \cuello{$+1$} \\
Bacha y grifería       & 4 & 3 & 1 & 0 \\
Lavavajillas + cocinas & 5 & 3 & 2 & 0 \\
\bottomrule
\end{tabular}
\end{table}

La relajación resulta exacta en este caso: como la cota de las mesadas se
obtiene sumando restricciones lineales, el problema relajado no puede superarla,
y la solución hallada ya es entera. El precio sombra del espacio de mesadas
—un combo por lugar— coincide con el resultado de la re-resolución de la
Tabla~\ref{tab:cap_e1}.

Esa coincidencia, sin embargo, es un resultado y no una garantía. En un modelo
entero el dual de la relajación puede sobreestimar el valor de un recurso, y de
hecho en el escenario siguiente el mismo lugar de mesada pasa a valer
\textbf{tres combos} en lugar de uno: ningún precio sombra de la relajación
anticipa ese salto. Es la razón por la cual el análisis de capacidad se realiza
re-resolviendo el modelo en cada escenario.

\subsection{Soluciones óptimas alternativas}
\label{sec:esc1_alt}

El modelo presenta degeneración: existen múltiples planes de stock que alcanzan
la misma variedad máxima. Se los determinó por enumeración exhaustiva,
evaluando las $\binom{20}{3} = 1.140$ ternas posibles contra las capacidades del
depósito, y también las $\binom{20}{4} = 4.845$ cuaternas.

\[
\underbrace{1.140}_{\text{ternas posibles}} \;-\;
\underbrace{\binom{13}{3} = 286}_{\text{ternas tomadas entre los combos 1 y 13}}
\;=\; \boxed{854 \text{ ternas óptimas}}
\qquad\qquad
0 \text{ cuaternas factibles}
\]

Las 286 ternas descartadas son las formadas exclusivamente por combos del 1 al
13, que ocupan seis lugares de lavavajillas y cocinas cuando el espacio admite
cinco.

\begin{table}[H]
\centering\small
\caption{Distribución de las 854 ternas óptimas.}
\label{tab:distribucion_e1}
\begin{tabular}{rr@{\qquad\qquad}rr@{\qquad\qquad}rr}
\toprule
\textbf{Unidades} & \textbf{Ternas} & \textbf{Lugares lavav.+coc.} & \textbf{Ternas} & \textbf{Categorías} & \textbf{Ternas} \\
\midrule
21 & 84  & 3 & 35  & 7 & 35 \\
22 & 385 & 4 & 273 & 8 & 819 \\
23 & 385 & 5 & 546 &   & \\
\bottomrule
\end{tabular}
\end{table}

Aplicando el criterio de desempate completo —ocho categorías cubiertas y 21
unidades— permanecen \textbf{49 ternas equivalentes}, de las cuales la
presentada es una. Las 35 ternas que cubren solo siete categorías coinciden
exactamente con las que ocupan tres lugares de lavavajillas y cocinas: son las
formadas únicamente por combos del 14 al 20, ninguno de los cuales incluye
lavavajillas.

Esto tiene dos consecuencias. La primera es de presentación: el plan de stock
debe ofrecerse como una solución representativa y no como la única posible. La
segunda es favorable para la operación: si la sección dispone de criterios que
el modelo no incorpora —márgenes, rotación o acuerdos con proveedores—, cuenta
con 49 planes alternativos entre los cuales elegir sin resignar variedad.

La enumeración cumple además una función de verificación: no utiliza el solver
sino que construye el stock requerido por cada conjunto de combos y lo compara
contra las capacidades. La coincidencia entre ambos procedimientos respalda la
formulación y no únicamente la resolución.

% ============================================================================
\section{Escenario 2 — Reposición automática}
\label{sec:esc2}

El nuevo contrato con proveedores de la zona permite considerar la reposición
como instantánea: lo que se vende se repone de inmediato. Esa modificación no
altera la estructura del modelo sino el supuesto S1. Si un artículo se repone al
instante, deja de ser necesario reservar un juego por combo y basta con
\textbf{una unidad de cada artículo}. En la formulación, el cambio afecta
únicamente a R1:

\[
\underbrace{x_p \geq \sum_c u_p\, a_{pc}\, y_c \quad \forall p}_{\text{reposición mensual}}
\qquad\Longrightarrow\qquad
\underbrace{x_p \geq u_p\, a_{pc}\, y_c \quad \forall p, \forall c}_{\text{reposición automática}}
\]

La cobertura pasa de una restricción por artículo (30) a una por par
artículo--combo (151), y el término de la derecha deja de acumularse. Un
artículo utilizado por diez combos ocupa un lugar y no diez.

\subsection{Variedad alcanzada y plan de stock}

La variedad máxima asciende a \textbf{13 combos sobre 20}, el 65\,\% del
catálogo. Quedan habilitados los combos 3, 4, 5, 7, 9, 10, 11, 12, 13, 14, 15,
16 y 17, con las ocho categorías cubiertas.

\begin{table}[H]
\centering\small
\caption{Efecto de la reposición automática sobre la oferta.}
\label{tab:comparacion_12}
\begin{tabular}{lrrr}
\toprule
\textbf{Indicador} & \textbf{Escenario 1} & \textbf{Escenario 2} & $\boldsymbol{\Delta}$ \\
\midrule
Combos disponibles para la venta   & 3  & 13 & $+10$ \\
Artículos distintos en el depósito & 17 & 27 & $+10$ \\
Unidades totales en stock          & 21 & 27 & $+6$ \\
Categorías del catálogo cubiertas  & 8  & 8  & 0 \\
Soluciones óptimas alternativas    & 854 & 4 & $-850$ \\
\bottomrule
\end{tabular}
\end{table}

\textbf{La variedad de oferta se ve beneficiada bajo cualquiera de sus dos
lecturas posibles.} Medida en combos disponibles, se multiplica por 4,3. Medida
en artículos distintos en el depósito, pasa de 17 a \textbf{27 de los 30},
quedando fuera únicamente \art{A4}, \art{M4} y \art{C4}. Bajo reposición
automática cada artículo utilizado requiere exactamente una unidad, de modo que
el plan de stock consiste en una unidad de cada uno de esos 27 artículos y las
unidades totales coinciden con la cantidad de artículos distintos.

El stock total se incrementa de 21 a 27 unidades, con lo cual la sección ocupa
más lugar del almacén compartido. La contrapartida es una mejora sustancial de
la eficiencia: se pasa de 7 unidades almacenadas por combo ofrecido a 2,1, al
eliminarse las unidades duplicadas de un mismo artículo.

\subsection{Ocupación del depósito}

\begin{table}[H]
\centering\small
\caption{Ocupación del depósito bajo reposición automática.}
\label{tab:ocupacion_e2}
\begin{tabular}{lrrrr}
\toprule
\textbf{Espacio} & \textbf{Capacidad} & \textbf{Ocupado} & \textbf{Holgura} & \textbf{Variantes existentes} \\
\midrule
Baldosas                  & 5 & 4 & 1 & 4 \\
Empapelado vinílico       & 8 & 4 & 4 & 4 \\
Apliques de luz           & 4 & 4 & 0 & 4 \\
Alacenas                  & 4 & 3 & 1 & 4 \\
\cuello{Mesadas}          & \cuello{3} & \cuello{3} & \cuello{0} & \cuello{4} \\
Bacha y grifería          & 4 & 4 & 0 & 4 \\
\cuello{Lavav. + cocinas} & \cuello{5} & \cuello{5} & \cuello{0} & \cuello{6} \\
\midrule
\textbf{Total} & \textbf{33} & \textbf{27} & \textbf{6} & \\
\bottomrule
\end{tabular}
\end{table}

La última columna de la Tabla~\ref{tab:ocupacion_e2} resulta indispensable para
interpretar correctamente las holguras, porque bajo reposición automática
\textbf{la capacidad de un espacio deja de acotar una cantidad de combos y pasa
a acotar una cantidad de variantes distintas} de esa categoría: cada variante
ocupa un lugar una sola vez, sea utilizada por un combo o por diez.

Bajo ese criterio, los cuatro espacios que aparecen completos no son
equivalentes. Apliques de luz y bacha y grifería disponen de exactamente cuatro
lugares para las cuatro variantes que existen, de modo que \textbf{no pueden
quedarse sin lugar}: su holgura nula únicamente indica que la solución utiliza
las cuatro variantes. La propia tabla lo confirma, porque alacenas presenta la
misma relación —cuatro lugares para cuatro variantes— y conserva una unidad de
holgura, ya que la solución prescinde de \art{A4}. \textbf{Ninguno de esos tres
espacios puede excluir un combo por sí solo.} Los cuellos de botella efectivos
son dos: \textbf{mesadas}, con tres lugares para cuatro variantes, y
\textbf{lavavajillas y cocinas}, con cinco lugares para seis variantes.

\subsection{Valor de una ampliación de capacidad}

\begin{table}[H]
\centering\small
\caption{Efecto de un lugar adicional bajo reposición automática.}
\label{tab:cap_e2}
\begin{tabular}{lrrr}
\toprule
\textbf{Espacio} & \textbf{Capacidad ampliada} & $\boldsymbol{V^{*}}$ & $\boldsymbol{\Delta V^{*}}$ \\
\midrule
Baldosas                  & 6 & 13 & 0 \\
Empapelado vinílico       & 9 & 13 & 0 \\
Apliques de luz           & 5 & 13 & 0 \\
Alacenas                  & 5 & 13 & 0 \\
\cuello{Mesadas}          & \cuello{4} & \cuello{16} & \cuello{$+3$} \\
Bacha y grifería          & 5 & 13 & 0 \\
\cuello{Lavav. + cocinas} & \cuello{6} & \cuello{15} & \cuello{$+2$} \\
\bottomrule
\end{tabular}
\end{table}

El resultado confirma el análisis anterior: únicamente mesadas y el espacio
compartido de lavavajillas y cocinas incrementan la variedad. Además lo hacen
\textbf{de a varios combos por lugar}, porque incorporar la variante faltante
habilita simultáneamente a todos los combos que la utilizan. Con la capacidad de
mesadas en cuatro lugares ese cuello de botella desaparece por completo, y el
único que subsiste es el de lavavajillas y cocinas.

La cantidad de soluciones óptimas alternativas se reduce de 854 a
\textbf{cuatro}. Bajo reposición automática el margen para almacenar unidades
sin utilizar es mucho menor, de modo que existen muchas menos formas de alcanzar
la misma variedad.

% ============================================================================
\section{Escenario 3 — Espacio compartido para mesadas y alacenas}
\label{sec:esc3}

La nueva disposición física del almacén permitiría reubicar las mesadas y las
alacenas en un espacio común \cita{con la capacidad de almacenar hasta 12
unidades en total}, bajo la misma modalidad con que actualmente se almacenan los
lavavajillas y las cocinas. El escenario se evalúa sobre la situación de
reposición automática.

El cambio es de datos y no de formulación: dos espacios de la lista se
consolidan en uno.

\[
\underbrace{\sum_j x_{Mj} \leq 3 \quad\text{y}\quad \sum_j x_{Aj} \leq 4}_{\text{7 lugares en 2 espacios}}
\qquad\Longrightarrow\qquad
\underbrace{\sum_j x_{Mj} + \sum_j x_{Aj} \leq 12}_{\text{12 lugares en 1 espacio}}
\]

El depósito pasa de 33 a \textbf{38 lugares} distribuidos en seis espacios.

\subsection{Variedad alcanzada}

La variedad máxima asciende a \textbf{16 combos sobre 20}, el 80\,\% del
catálogo. Quedan habilitados los combos 1, 2, 3, 4, 5, 7, 9, 10, 11, 12, 13, 14,
15, 16, 17 y 18, con 29 unidades en stock: una unidad de 29 de los 30 artículos
del catálogo.

\begin{table}[H]
\centering\small
\caption{Efecto de la nueva disposición física sobre la oferta.}
\label{tab:comparacion_23}
\begin{tabular}{lrrr}
\toprule
\textbf{Indicador} & \textbf{Escenario 2} & \textbf{Escenario 3} & $\boldsymbol{\Delta}$ \\
\midrule
Combos disponibles para la venta  & 13 & 16 & $+3$ \\
Unidades totales en stock         & 27 & 29 & $+2$ \\
Categorías del catálogo cubiertas & 8  & 8  & 0 \\
Soluciones óptimas alternativas   & 4  & 1  & $-3$ \\
\bottomrule
\end{tabular}
\end{table}

\textbf{La nueva disposición no permite completar la oferta de combos.} Los
cuatro que quedan excluidos son el \textbf{6, el 8, el 19 y el 20}, y la causa
es precisa: son exactamente los cuatro combos que incluyen \art{C4}, la cocina a
gas negra, único artículo del catálogo que permanece en cero. El espacio de
lavavajillas y cocinas conserva cinco lugares para seis variantes, de modo que
una de ellas debe resignarse, y la resignada es \art{C4} por ser la que menos
combos habilita: cuatro, frente a entre cinco y ocho de las demás variantes de
ese espacio (Tabla~\ref{tab:apariciones}).

\subsection{Ocupación del depósito}

\begin{table}[H]
\centering\small
\caption{Ocupación del depósito con mesadas y alacenas consolidadas.}
\label{tab:ocupacion_e3}
\begin{tabular}{lrrrr}
\toprule
\textbf{Espacio} & \textbf{Capacidad} & \textbf{Ocupado} & \textbf{Holgura} & \textbf{Variantes existentes} \\
\midrule
Baldosas                  & 5 & 4 & 1 & 4 \\
Empapelado vinílico       & 8 & 4 & 4 & 4 \\
Apliques de luz           & 4 & 4 & 0 & 4 \\
Bacha y grifería          & 4 & 4 & 0 & 4 \\
\cuello{Lavav. + cocinas} & \cuello{5} & \cuello{5} & \cuello{0} & \cuello{6} \\
Mesadas + alacenas        & 12 & 8 & 4 & 8 \\
\midrule
\textbf{Total} & \textbf{38} & \textbf{29} & \textbf{9} & \\
\bottomrule
\end{tabular}
\end{table}

El espacio consolidado dispone de 12 lugares para un máximo de 8 variantes
—cuatro de mesada y cuatro de alacena—, de modo que \textbf{deja de restringir
el problema en cualquier configuración}, y conserva cuatro lugares libres en la
solución final. Apliques de luz y bacha y grifería vuelven a presentar holgura
nula sin que ello implique restricción alguna: su capacidad iguala el número de
variantes existentes, de modo que no pueden quedarse sin lugar. El único cuello
de botella efectivo que subsiste es el espacio de lavavajillas y cocinas.

\subsection{Conveniencia de la reubicación}

\textbf{La nueva disposición física es conveniente.} La variedad se incrementa
de 13 a 16 combos —un 23\,\%— \textbf{sin requerir capacidad adicional}: los
siete lugares hoy repartidos entre mesadas (3) y alacenas (4) pasan a ser doce
en un espacio único, con lo cual ambas categorías dejan de limitar la oferta. La
reubicación no demanda mercadería ni superficie adicional, y su rendimiento es
de tres combos.

\subsection{Capacidad adicional a solicitar}

\begin{table}[H]
\centering\small
\caption{Efecto de un lugar adicional bajo la nueva disposición.}
\label{tab:cap_e3}
\begin{tabular}{lrrr}
\toprule
\textbf{Espacio} & \textbf{Capacidad ampliada} & $\boldsymbol{V^{*}}$ & $\boldsymbol{\Delta V^{*}}$ \\
\midrule
Baldosas                  & 6  & 16 & 0 \\
Empapelado vinílico       & 9  & 16 & 0 \\
Apliques de luz           & 5  & 16 & 0 \\
Bacha y grifería          & 5  & 16 & 0 \\
\cuello{Lavav. + cocinas} & \cuello{6} & \cuello{20} & \cuello{$+4$} \\
Mesadas + alacenas        & 13 & 16 & 0 \\
\bottomrule
\end{tabular}
\end{table}

\textbf{Un único lugar adicional en el espacio de lavavajillas y cocinas —de
cinco a seis— es suficiente para completar los 20 combos del catálogo.} La
ampliación recupera de manera simultánea los cuatro combos que incluyen
\art{C4}, dado que con seis lugares las seis variantes de esa categoría pueden
almacenarse en conjunto.

Ningún otro espacio requiere ampliación: el espacio consolidado de mesadas y
alacenas conserva cuatro lugares libres, y a los restantes puede sumárseles
capacidad sin que la variedad se modifique. Solicitar cualquier otra unidad de
almacenamiento no produciría ningún efecto sobre la oferta.

Bajo esta disposición la solución óptima además es \textbf{única}: de las 4.845
combinaciones posibles de 16 combos, una sola resulta factible. Al subsistir una
única restricción activa en todo el modelo, la única forma de alcanzar 16 combos
consiste en tomar los 16 que no utilizan \art{C4}.

% ============================================================================
\section{Sensibilidad del modelo a los supuestos}
\label{sec:sensibilidad}

Los supuestos que admiten más de una lectura se evaluaron resolviendo el modelo
completo bajo cada alternativa. El supuesto S1 —la interpretación de la
disponibilidad— ya quedó cuantificado por la comparación entre los
Escenarios~\ref{sec:esc1} y~\ref{sec:esc2}, que constituye en sí misma su
prueba: 3 combos frente a 13. Esta sección aborda los dos restantes.

\subsection{S2 — Unidades de empapelado por combo}
\label{sec:sens_s2}

\begin{table}[H]
\centering\small
\caption{Variedad máxima según las unidades de empapelado que requiere un combo.}
\label{tab:sens_s2}
\begin{tabular}{rrlrr}
\toprule
\textbf{Rollos por combo} & $\boldsymbol{V^{*}}$ & \textbf{Combos habilitados} & \textbf{Con empapelado} & \textbf{Categorías} \\
\midrule
1 (base) & 3 & 12 -- 15 -- 18 & 2 & 8 \\
2 a 8    & 3 & 12 -- 13 -- 15 & 1 & 8 \\
9        & 2 & 12 -- 13       & 0 & 7 \\
\bottomrule
\end{tabular}
\end{table}

La variedad no se modifica hasta ocho rollos por combo, un valor ya fuera de
todo rango plausible. El mecanismo que explica esa insensibilidad es la
composición del catálogo: los combos 12 y 13 no incluyen empapelado y completan
la terna sin ocupar ese espacio. \textbf{El supuesto S2 resulta irrelevante para
el resultado en cualquier rango razonable.}

La última fila aporta además la justificación empírica de por qué la cobertura
de categorías se incorporó como objetivo y no como restricción: con nueve rollos
por combo, ninguna solución factible puede cubrir el empapelado. Formulada como
restricción, el modelo resultaría infactible y no produciría información alguna;
como etapa de desempate, informa que cubre siete categorías y continúa
resolviendo.

\subsection{S3 — Criterio de desempate}
\label{sec:sens_s3}

\begin{table}[H]
\centering\small
\caption{Resultado de los cuatro criterios de desempate evaluados, sobre el
escenario base.}
\label{tab:sens_s3}
\begin{tabular}{llrrr}
\toprule
\textbf{Criterio} & \textbf{Combos} & $\boldsymbol{V^{*}}$ & \textbf{Categorías} & \textbf{Unidades} \\
\midrule
Categorías y unidades (adoptado) & 12 -- 15 -- 18 & 3 & 8 & 21 \\
Mínimo de unidades               & 14 -- 18 -- 20 & 3 & 7 & 21 \\
Máximo de artículos              & \ 5 -- \ 9 -- 15 & 3 & 8 & 23 \\
Sin desempate                    & 15 -- 17 -- 20 & 3 & 7 & 21 \\
\bottomrule
\end{tabular}
\end{table}

Los cuatro criterios alcanzan la misma variedad máxima: el desempate nunca le
cuesta un combo a la sección, únicamente determina cuál de las soluciones
equivalentes se adopta. El criterio elegido es el único que garantiza la
cobertura de las ocho categorías al mínimo costo de espacio; el de máximo de
artículos también las cubre, pero ocupando dos lugares adicionales.

Cabe señalar que cubrir las ocho categorías \textbf{no representa costo alguno
en términos de depósito}: tanto el criterio adoptado como el de mínimo de
unidades ocupan 21 lugares. Lo único que se modifica es la distribución interna
—un lugar más en lavavajillas y cocinas y uno menos en empapelado—.

La última fila merece una advertencia. Sin criterio de desempate, la solución
obtenida resultó tener 21 unidades por coincidencia; nada en el modelo lo
garantiza, y otra versión del solver podría devolver cualquiera de las 854
ternas óptimas, incluida alguna de 23 unidades.

% ============================================================================
\section{Conclusiones y recomendaciones}
\label{sec:conclusiones}

\begin{table}[H]
\centering\small
\caption{Síntesis de los tres escenarios analizados.}
\label{tab:sintesis}
\begin{tabular}{llrrr}
\toprule
\textbf{Escenario} & \textbf{Espacio que limita} & \textbf{Combos} & \textbf{\% catálogo} & \textbf{Óptimos} \\
\midrule
1 — Reposición mensual        & Mesadas                        & 3  & 15\,\%  & 854 \\
2 — Reposición automática     & Mesadas y lavav. + cocinas     & 13 & 65\,\%  & 4 \\
3 — Mesadas y alacenas juntas & Lavavajillas + cocinas         & 16 & 80\,\%  & 1 \\
\midrule
3 con un lugar adicional      & ---                            & 20 & 100\,\% & --- \\
\bottomrule
\end{tabular}
\end{table}

Los tres escenarios convergen en un mismo diagnóstico: \textbf{la Sección
Cocinas no enfrenta una limitación de capacidad general sino la de un único
espacio por vez}. En cada configuración existe un solo cuello de botella que
determina íntegramente el resultado, mientras el resto del almacén conserva
lugar sin utilizar —12 de 33 lugares en el escenario base—. Esa observación
cambia la naturaleza de cualquier negociación de espacio con las demás secciones
de la sucursal: lo relevante no es cuánto lugar se solicita sino en qué espacio.

\textbf{La decisión de modelado de mayor incidencia no fue de cálculo sino de
interpretación.} Entender la disponibilidad bajo reposición mensual como reserva
exclusiva de unidades, y no como una unidad compartida entre todos los combos
que la utilizan, es lo que separa 3 combos de 13. Ningún dato del problema
resuelve esa lectura: se adoptó con justificación explícita
(Sección~\ref{sec:supuestos}) y se cuantificó su efecto mediante la comparación
de escenarios.

\textbf{El objetivo de variedad, por sí solo, no determina la decisión
buscada.} Maximizar la cantidad de combos define correctamente cuántos pueden
ofrecerse, pero no cuántas unidades mantener en stock, que es la decisión
operativa. Resuelto únicamente ese funcional, el modelo ocupa el depósito
completo con 12 unidades que no habilitan ningún combo. Las etapas
lexicográficas incorporadas cierran esa indeterminación, y la de cobertura de
categorías lo hace sin costo alguno en espacio.

\textbf{La ausencia de análisis post-óptimo no restringió el análisis de
sensibilidad.} El procedimiento de re-resolución resultó, de hecho, más
informativo que los precios sombra: un lugar de mesada vale un combo en el
escenario base y tres en el de reposición automática, diferencia que el dual de
la relajación lineal no habría anticipado. El costo es computacional —una
resolución por escenario— y a esta escala resulta despreciable.

\paragraph{Recomendaciones.} Se desprenden tres cursos de acción, en este orden
de prioridad:

\begin{enumerate}
  \item \textbf{Operar bajo el nuevo esquema de reposición automática.} Eleva la
  oferta de 3 a 13 combos sin modificar la disposición del depósito ni solicitar
  capacidad adicional.
  \item \textbf{Ejecutar la reubicación de mesadas y alacenas} en el espacio
  compartido de 12 unidades. Suma tres combos más, también sin requerir
  capacidad adicional.
  \item \textbf{Solicitar un lugar adicional en el espacio de lavavajillas y
  cocinas.} Es la única ampliación con efecto sobre la oferta, y completa los 20
  combos del catálogo.
\end{enumerate}

% ============================================================================
\section*{Anexo — Herramientas y verificación}
\addcontentsline{toc}{section}{Anexo — Herramientas y verificación}

\paragraph{Entorno y herramientas.} El trabajo se desarrolló en Python 3.10 con
el siguiente conjunto de librerías:

\begin{itemize}
  \item \textbf{PuLP} \cite{pulp} con el solver \textbf{CBC} \cite{cbc}, para el
  modelado y la resolución del problema de programación lineal entera. CBC
  implementa internamente el algoritmo de \emph{Branch \& Bound}. La
  correspondencia con la notación de LINDO/LINGO es directa: \texttt{INT y} se
  expresa como \texttt{LpVariable(..., cat=LpBinary)}, \texttt{GIN x} como
  \texttt{cat=LpInteger}, y los precios duales se obtienen mediante
  \texttt{constraint.pi}, con validez restringida a la relajación lineal.
  \item \textbf{pandas}, para la construcción de las tablas de resultados.
  \item \textbf{matplotlib}, para el material gráfico de trabajo.
  \item \texttt{itertools.combinations}, de la biblioteca estándar, para la
  enumeración exhaustiva de soluciones óptimas alternativas.
\end{itemize}

\paragraph{Diseño del modelo.} La formulación se escribió una única vez, como
función parametrizada:
\begin{quote}
\centering\small
\texttt{resolver(params, objetivo, variedad\_min, categorias\_min, combos\_fijos, relajar)}
\end{quote}
Las tres etapas lexicográficas, la relajación lineal, los análisis de capacidad,
las pruebas de supuestos y los tres escenarios son llamadas a esa misma función
con distintos parámetros. Los supuestos y las capacidades residen en un único
diccionario de configuración, y la solicitud de un supuesto, un espacio o una
categoría inexistente produce un error en lugar de devolver silenciosamente el
caso base. De los tres escenarios, únicamente el tercero requirió modificar el
código, y en una sola línea: permitir que la configuración reciba la estructura
completa de espacios y no solo sus capacidades. El resto del programa
—restricciones, cotas, ocupación y reportes— recorre la lista de espacios de
forma genérica.

\paragraph{Verificación.} Los resultados se validaron por tres vías
independientes entre sí. La primera es \textbf{analítica}: la cota
$\sum_c y_c \leq 3$ deducida en la Sección~\ref{sec:restricciones} predice el
óptimo del escenario base sin intervención del solver. La segunda es la
\textbf{enumeración exhaustiva} de la Sección~\ref{sec:esc1_alt}, que tampoco lo
utiliza. La tercera es un conjunto de \textbf{controles automáticos} que se
ejecutan antes de aceptar cualquier resultado:

\begin{itemize}
  \item \textbf{46 sobre los datos de entrada}: que todo artículo integrante de
  un combo exista en el catálogo, que ningún combo repita categoría, las 30
  frecuencias de la Tabla~\ref{tab:apariciones} una por una, las capacidades y
  las cotas de la Tabla~\ref{tab:cota_trivial}.
  \item \textbf{25 cruzados sobre los óptimos} de los tres escenarios: que la
  variedad no supere la cota analítica, que el stock de la solución final sea
  exactamente el necesario, que la ocupación respete todas las capacidades y que
  la solución figure entre las enumeradas exhaustivamente.
  \item \textbf{10 sobre las pruebas de supuestos}, formulados como propiedades
  que cualquier resultado correcto debe satisfacer: ampliar la capacidad nunca
  puede reducir la variedad, incrementar las unidades por combo nunca puede
  aumentarla, ningún criterio de desempate puede perder variedad.
\end{itemize}

Estos últimos no verifican que el resultado coincida con el esperado sino
propiedades que todo resultado correcto debe cumplir, lo que los vuelve útiles
precisamente cuando el resultado sorprende.

% ============================================================================
\renewcommand{\refname}{Bibliografía}
\addcontentsline{toc}{section}{Bibliografía}

\begin{thebibliography}{9}

\bibitem{hillier} Hillier, F. S. y Lieberman, G. J. (2015).
  \emph{Introducción a la Investigación de Operaciones} (9.\textsuperscript{a}
  ed.). México: McGraw-Hill Interamericana. Capítulo 11: programación entera,
  formulación con variables binarias, algoritmo de \emph{Branch \& Bound} y
  empleo de la relajación lineal como cota.

\bibitem{catedra-metas} Cátedra de Optimización, UCA (2026).
  \emph{Clase 5 — Programación por metas}. Material de cátedra,
  2.\textsuperscript{o} cuatrimestre de 2026. Fuente del esquema de metas
  lexicográficas aplicado en la Sección~\ref{sec:objetivo}.

\bibitem{catedra-mip} Cátedra de Optimización, UCA (2026).
  \emph{Clase 6 — Programación entera y binaria}. Material de cátedra,
  2.\textsuperscript{o} cuatrimestre de 2026. Fuente del algoritmo de resolución
  y de la inaplicabilidad del análisis post-óptimo en modelos enteros.

\bibitem{pulp} Mitchell, S., O'Sullivan, M. y Dunning, I.
  \emph{PuLP: A Linear Programming Toolkit for Python}. Documentación oficial
  del paquete \texttt{PuLP}, utilizado como motor de modelado y resolución.

\bibitem{cbc} Forrest, J. y Lougee-Heimer, R.
  \emph{CBC (COIN-OR Branch and Cut)}. Solver de programación lineal y entera
  empleado por PuLP para la resolución de todos los modelos de este informe.

\end{thebibliography}

\end{document}
